%=============================================================================
% ABSTRACT -- background, gap, proposal, key results, code link (~200 words).
%=============================================================================
\begin{abstract}
District-level dengue forecasts three weeks ahead give Sri Lankan health services
time to act. Spatio-temporal graph neural networks (ST-GNNs) and SEIR-informed
models have been proposed for this task, but they predict the case level directly
and are evaluated without the persistence forecast, even though last week's count
explains $r^2=0.89$ of the next. We propose \model{} (Persistence-Anchored Gated
Estimation), a forecasting head for any spatio-temporal encoder that predicts a
gated departure from the last observed week, at the cost of one parameter. We
correct the benchmark into a leakage-free 559-week dataset, extend it to 2026 by
parsing 127 new reports, and
evaluate five published ST-GNNs and an LSTM under one protocol. \model{} repairs
the three encoders that fail as published by 5.95--16.29 validation RMSE (9/9 paired
runs each), and with AAGCN it beats persistence on all nine runs on validation and
test. An ablation attributes the gain to the anchor and the gate; an SEIR simulator
in the same head adds nothing, and on 85 later forecast starts it ties its no-SEIR
twin ($p=0.67$). On those 2024--2026 weeks no learned model beats persistence on
RMSE (34.75 against 35.43 for \model{}), although \model{} has lower MAE and SMAPE.
For this target, persistence is the prior that matters. Code and data: \codeurl.
\end{abstract}

\begin{IEEEkeywords}
dengue forecasting, spatio-temporal graph neural networks, persistence baseline,
physics-informed learning, data leakage, prospective evaluation
\end{IEEEkeywords}
